
\subsection{Carry a loss quantity without inventing a tax asset}

The negative expression of $-19$ needs a rule before it can affect a later
tax calculation. The finite extension supplies a conditional loss ledger.
Each lot has an identity, an origin reference, an availability date, an explicit
expiry declaration, original and remaining quantities, cumulative usage and
expiry, and an eligibility finding. An unknown finding remains unresolved;
a false finding is unsupported. A generated loss therefore becomes a quantity
in this separate ledger only under an affirmative, traceable assumption.
It is not a cash refund or a recognized deferred-tax asset.

The selected illustrative policy evaluates all declared usage on the last
included calendar day of an interval. An exclusive expiry on that day makes
the lot unavailable for usage. Expiry at the following endpoint permits usage
on the preceding day, but the remainder expires before the next opening.
No expiry is itself an explicit assumption. These conventions specify this
finite model; another jurisdiction or an intraperiod method may require a
different rule. Neither FIFO nor a legal utilization percentage is assumed.
The author supplies the chosen lot allocations and an interval usage cap.

For each lot, original loss equals remaining, used and expired quantities.
Across a step, with $L_0$ and $L_1$ denoting opening and closing remaining
quantities, $G$ generation, $U$ usage and $E$ expiry,
\[
 L_0+G=U+E+L_1.
\]
All quantities on this line are nonnegative exact rationals. A generated loss
becomes available at the exclusive endpoint and cannot be used in its own
generation interval. A usage cannot exceed its lot's remaining quantity,
the declared interval cap or positive taxable profit. Zero and exhausted lots
remain identified, preventing the same identity from being introduced again.
Consecutive ledger entries must agree on both interval endpoints and their
opening and closing quantities.

The checked negative example generates 19 under explicit eligibility and
no-expiry premises. A separate hypothetical 2026 calculation supplies taxable
profit of 21. Explicit use of 19 leaves taxable profit of 2 and gross conditional
tax of $1/2$ at rate $1/4$. This connects the loss ledger across periods; it does
not establish financial-book carry. The 2026 financial opening is independently
declared for the test and is expressly not the carried closing state of the
2024--2026 example. Both complete conditional calculations remain attached.

Partial use makes the expiry rule consequential. Using 7 of 19 leaves 12 when
the expiry is later than the next endpoint. If the exclusive expiry equals that
endpoint, 7 may be used on the preceding day and the remaining 12 expires.
If expiry instead equals the usage day, the usage is rejected; with no usage,
all 19 expires. These outcomes follow the selected policy, rather than a
silent decision to use the oldest lot or to preserve an expired balance.

The operation proposes a complete next loss state and per-lot movements.
Applying it recomputes the proposal and compares the complete result, including
the original financial calculation. Changed amounts, origins, eligibility,
scope, cursors or even Boolean-to-integer substitutions are rejected. This is
a pure in-memory application: it creates no durable commit, tax journal,
deferred asset or annual authorization. Connecting the conditional tax amount
and loss ledger to an actual financial tax balance still requires separately
specified recognition, accrual and settlement methods.
